\section{Lecture 25: Variational Autoencoders}

\subsection{Motivation: What Is Generative Modeling?}

So far, much of machine learning has focused on prediction:
\[
x \mapsto y.
\]

In generative modeling, the goal is different.

\medskip

\noindent
\textbf{Goal.}  
Learn the structure of the data distribution itself.

\medskip

\noindent
Instead of only predicting labels, we want a model that captures
\[
p_{\text{data}}(x),
\]
or an approximation
\[
p_\theta(x) \approx p_{\text{data}}(x).
\]

\medskip

\noindent
\textbf{What this enables.}

\begin{itemize}
    \item sampling new data points
    \item modeling uncertainty
    \item learning useful latent representations
    \item interpolation and structured generation
\end{itemize}

\medskip

\noindent
\textbf{Examples.}

\begin{itemize}
    \item generate images similar to a training dataset
    \item compress data into meaningful hidden factors
    \item fill in missing observations
\end{itemize}

\medskip

\noindent
\textbf{Key idea.}  
A generative model does not merely separate classes; it tries to explain how data could have been produced.

\subsection{Latent-Variable Viewpoint}

A powerful way to model complicated data is to assume that observations are generated from hidden variables.

\medskip

\noindent
\textbf{Latent-variable model.}

\[
p_\theta(x,z) = p_\theta(x|z)p(z),
\]
where:
\begin{itemize}
    \item $z$ is a latent variable
    \item $p(z)$ is a prior distribution
    \item $p_\theta(x|z)$ is the conditional data model
\end{itemize}

\medskip

\noindent
The observable distribution is obtained by marginalizing out $z$:
\[
p_\theta(x) = \int p_\theta(x,z)\,dz
= \int p_\theta(x|z)p(z)\,dz.
\]

\medskip

\noindent
\textbf{Interpretation.}  
The latent variable $z$ captures underlying factors of variation, while the conditional model $p_\theta(x|z)$ turns those factors into an actual observation.

\medskip

\noindent
\textbf{Example intuition.}

In image generation, $z$ might encode:
\begin{itemize}
    \item pose
    \item shape
    \item style
    \item lighting
\end{itemize}

and the decoder maps these factors into an image $x$.

\medskip

\noindent
\textbf{Challenge.}  
Even if the model is conceptually simple, the marginal likelihood
\[
p_\theta(x)=\int p_\theta(x|z)p(z)\,dz
\]
is often hard to compute exactly.

\subsection{Likelihood and the Need for Approximate Inference}

If we want to learn $\theta$ by maximum likelihood, we would like to maximize
\[
\log p_\theta(x)
=
\log \int p_\theta(x,z)\,dz.
\]

\medskip

\noindent
\textbf{Difficulty.}  
The integral over $z$ is usually intractable when:
\begin{itemize}
    \item $z$ is high-dimensional
    \item $p_\theta(x|z)$ is parameterized by a neural network
    \item the posterior $p_\theta(z|x)$ has no closed form
\end{itemize}

\medskip

\noindent
This creates two coupled problems:

\begin{itemize}
    \item \textbf{inference:} estimate which latent variables explain a given $x$
    \item \textbf{learning:} update model parameters to better fit the data
\end{itemize}

\medskip

\noindent
\textbf{Bridge to variational inference.}  
Instead of computing the exact posterior
\[
p_\theta(z|x)=\frac{p_\theta(x,z)}{p_\theta(x)},
\]
we introduce a tractable approximation
\[
q_\phi(z|x).
\]

\medskip

\noindent
This leads directly to the central idea behind variational autoencoders.

\subsection{From Autoencoders to Variational Autoencoders}

A classical autoencoder consists of:
\begin{itemize}
    \item an encoder: $x \mapsto h$
    \item a decoder: $h \mapsto \hat{x}$
\end{itemize}

trained so that $\hat{x}$ reconstructs $x$.

\medskip

\noindent
\textbf{Limitation of a classical autoencoder.}  
It is usually a deterministic compression model, not a full probabilistic generative model.

\medskip

\noindent
\textbf{VAE idea.}  
Replace the deterministic hidden code by a probability distribution over latent variables:
\[
q_\phi(z|x).
\]

Then:
\begin{itemize}
    \item the \textbf{encoder} produces a distribution over latent codes
    \item the \textbf{decoder} defines a distribution over reconstructed data
\end{itemize}

\medskip

\noindent
Thus a VAE combines:
\begin{itemize}
    \item latent-variable generative modeling
    \item variational inference
    \item neural-network parameterization
\end{itemize}

\subsection{The Generative Model and the Approximate Posterior}

In a VAE, we specify:

\medskip

\noindent
\textbf{Prior.}
\[
p(z)=\mathcal{N}(0,I)
\]
in the standard setting.

\medskip

\noindent
\textbf{Decoder / generative model.}
\[
p_\theta(x|z).
\]

\medskip

\noindent
\textbf{Encoder / approximate posterior.}
\[
q_\phi(z|x).
\]

\medskip

\noindent
The decoder answers:
\begin{quote}
Given a latent code $z$, what distribution over data $x$ should we generate?
\end{quote}

The encoder answers:
\begin{quote}
Given an observed data point $x$, what latent values $z$ are likely to have produced it?
\end{quote}

\medskip

\noindent
\textbf{Typical Gaussian encoder.}
\[
q_\phi(z|x)=\mathcal{N}(z\mid \mu_\phi(x), \Sigma_\phi(x)).
\]

Often one uses diagonal covariance:
\[
\Sigma_\phi(x)=\mathrm{diag}(\sigma_\phi^2(x)).
\]

\medskip

\noindent
\textbf{Insight.}  
The encoder performs approximate inference, while the decoder performs generation.

\subsection{Deriving the ELBO}

We now derive the objective used to train a VAE.

\medskip

\noindent
Start from the marginal log-likelihood:
\[
\log p_\theta(x)=\log \int p_\theta(x,z)\,dz.
\]

Introduce any distribution $q_\phi(z|x)$ and rewrite:
\[
\log p_\theta(x)
=
\log \int q_\phi(z|x)\frac{p_\theta(x,z)}{q_\phi(z|x)}\,dz.
\]

Applying Jensen's inequality gives:
\[
\log p_\theta(x)
\ge
\int q_\phi(z|x)\log \frac{p_\theta(x,z)}{q_\phi(z|x)}\,dz.
\]

\medskip

\noindent
\textbf{Definition (ELBO).}
\[
\mathcal{L}(x;\theta,\phi)
=
\mathbb{E}_{q_\phi(z|x)}
\left[
\log \frac{p_\theta(x,z)}{q_\phi(z|x)}
\right].
\]

This is called the \textbf{Evidence Lower Bound}.

\medskip

\noindent
Using
\[
p_\theta(x,z)=p_\theta(x|z)p(z),
\]
we obtain
\[
\mathcal{L}(x;\theta,\phi)
=
\mathbb{E}_{q_\phi(z|x)}[\log p_\theta(x|z)]
+
\mathbb{E}_{q_\phi(z|x)}[\log p(z)]
-
\mathbb{E}_{q_\phi(z|x)}[\log q_\phi(z|x)].
\]

Grouping the last two terms,
\[
\mathcal{L}(x;\theta,\phi)
=
\mathbb{E}_{q_\phi(z|x)}[\log p_\theta(x|z)]
-
\mathrm{KL}\!\left(q_\phi(z|x)\,\|\,p(z)\right).
\]

\medskip

\noindent
\textbf{This is the standard VAE objective.}

\subsection{Why the ELBO Makes Sense}

The ELBO can also be related directly to the true posterior.

\medskip

\noindent
\textbf{Key decomposition.}
\[
\log p_\theta(x)
=
\mathcal{L}(x;\theta,\phi)
+
\mathrm{KL}\!\left(q_\phi(z|x)\,\|\,p_\theta(z|x)\right).
\]

Since KL divergence is always nonnegative,
\[
\mathrm{KL}\!\left(q_\phi(z|x)\,\|\,p_\theta(z|x)\right)\ge 0,
\]
we obtain
\[
\mathcal{L}(x;\theta,\phi)\le \log p_\theta(x).
\]

\medskip

\noindent
\textbf{Interpretation.}  
Maximizing the ELBO does two things at once:

\begin{itemize}
    \item it pushes the model toward higher data likelihood
    \item it makes the approximate posterior $q_\phi(z|x)$ closer to the true posterior $p_\theta(z|x)$
\end{itemize}

\medskip

\noindent
\textbf{Insight.}  
The ELBO is not an arbitrary training loss. It is a principled surrogate for maximum likelihood in a latent-variable model.

\subsection{Encoder--Decoder Structure}

The architecture of a VAE is conceptually simple.

\medskip

\noindent
\textbf{Encoder network.}

Given $x$, the encoder outputs parameters of a latent distribution:
\[
\mu_\phi(x), \quad \log \sigma_\phi^2(x).
\]

This defines
\[
q_\phi(z|x)=\mathcal{N}(z\mid \mu_\phi(x), \mathrm{diag}(\sigma_\phi^2(x))).
\]

\medskip

\noindent
\textbf{Sampling step.}

A latent sample $z$ is drawn from this distribution.

\medskip

\noindent
\textbf{Decoder network.}

Given $z$, the decoder outputs parameters of a distribution over $x$:
\[
p_\theta(x|z).
\]

For example:
\begin{itemize}
    \item Gaussian output for continuous data
    \item Bernoulli output for binary data
\end{itemize}

\medskip

\noindent
\textbf{Learning objective.}  
Jointly optimize $\theta$ and $\phi$ to maximize the ELBO over the dataset:
\[
\sum_{i=1}^n \mathcal{L}(x_i;\theta,\phi).
\]

\medskip

\noindent
\textbf{Interpretation of the architecture.}

\begin{itemize}
    \item encoder: approximate inference
    \item decoder: probabilistic generation
    \item latent space: compressed structured representation
\end{itemize}

\subsection{Reconstruction Term and KL Term}

The ELBO has two parts:
\[
\mathcal{L}(x;\theta,\phi)
=
\underbrace{\mathbb{E}_{q_\phi(z|x)}[\log p_\theta(x|z)]}_{\text{reconstruction term}}
-
\underbrace{\mathrm{KL}(q_\phi(z|x)\|p(z))}_{\text{regularization term}}.
\]

\medskip

\noindent
\textbf{Reconstruction term.}  
This encourages the decoder to assign high probability to the observed data when conditioned on latent codes inferred from $x$.

\medskip

\noindent
\textbf{KL term.}  
This encourages the encoder's posterior distribution to remain close to the prior.

\medskip

\noindent
\textbf{Tradeoff.}

\begin{itemize}
    \item strong emphasis on reconstruction $\rightarrow$ better fidelity, weaker latent structure
    \item strong emphasis on KL regularization $\rightarrow$ smoother latent space, but risk of underusing the latent code
\end{itemize}

\medskip

\noindent
\textbf{Insight.}  
A VAE is not just trying to reconstruct data. It is trying to reconstruct data using a latent code that remains organized and compatible with a simple prior.

\subsection{The Reparameterization Trick}

A challenge arises because the ELBO contains an expectation over
\[
z \sim q_\phi(z|x),
\]
and $q_\phi$ depends on the encoder parameters $\phi$.

\medskip

\noindent
\textbf{Problem.}  
Naively sampling $z$ can obstruct gradient flow through the encoder.

\medskip

\noindent
\textbf{Solution.}  
Rewrite the random variable so that the randomness is separated from the parameters.

For a Gaussian encoder,
\[
z \sim \mathcal{N}(\mu,\sigma^2 I)
\]
can be written as
\[
z = \mu + \sigma \odot \epsilon,
\qquad
\epsilon \sim \mathcal{N}(0,I).
\]

Thus
\[
z = \mu_\phi(x) + \sigma_\phi(x)\odot \epsilon.
\]

\medskip

\noindent
\textbf{Equivalent expectation.}
\[
\mathbb{E}_{q_\phi(z|x)}[f(z)]
=
\mathbb{E}_{\epsilon \sim \mathcal{N}(0,I)}
\left[
f\bigl(\mu_\phi(x)+\sigma_\phi(x)\odot \epsilon\bigr)
\right].
\]

\medskip

\noindent
\textbf{Why this helps.}  
Now the stochasticity is entirely in $\epsilon$, which does not depend on $\phi$. Therefore gradients can backpropagate through $\mu_\phi(x)$ and $\sigma_\phi(x)$.

\medskip

\noindent
\textbf{Insight.}  
The reparameterization trick turns sampling into a differentiable computation graph.

\subsection{Closed-Form KL for a Gaussian Encoder}

In the standard VAE setting,
\[
q_\phi(z|x)=\mathcal{N}(\mu,\mathrm{diag}(\sigma^2)),
\qquad
p(z)=\mathcal{N}(0,I).
\]

Then the KL divergence has a closed form:
\[
\mathrm{KL}\!\left(q_\phi(z|x)\,\|\,p(z)\right)
=
\frac{1}{2}\sum_{j=1}^d
\left(
\mu_j^2+\sigma_j^2-\log \sigma_j^2 -1
\right).
\]

\medskip

\noindent
\textbf{Practical significance.}

\begin{itemize}
    \item no Monte Carlo estimator is needed for the KL term
    \item optimization becomes efficient and stable
\end{itemize}

\medskip

\noindent
Typically the reconstruction term is estimated with one or a few samples of $\epsilon$, while the KL term is computed exactly.

\subsection{How to Read the KL Term}

The KL penalty is often described as a regularizer, but it has several connected meanings.

\medskip

\noindent
\textbf{First interpretation: closeness to the prior.}  
It prevents the encoder from producing arbitrary latent distributions far from $p(z)$.

\medskip

\noindent
\textbf{Second interpretation: structure in latent space.}  
If $q_\phi(z|x)$ stays near a shared simple prior, nearby latent points tend to decode into meaningful nearby outputs.

\medskip

\noindent
\textbf{Third interpretation: information bottleneck flavor.}  
A large KL term means the posterior is carrying substantial information about the specific input $x$. Penalizing this limits how much information can be stored in $z$.

\medskip

\noindent
\textbf{Practical consequence.}

\begin{itemize}
    \item too little KL pressure $\rightarrow$ latent space becomes irregular, poor generative sampling
    \item too much KL pressure $\rightarrow$ latent code carries too little information
\end{itemize}

\subsection{Training and Generation}

\textbf{Training phase.}

For each data point $x$:
\begin{enumerate}
    \item compute encoder outputs $\mu_\phi(x)$ and $\sigma_\phi(x)$
    \item sample $\epsilon \sim \mathcal{N}(0,I)$
    \item construct
    \[
    z=\mu_\phi(x)+\sigma_\phi(x)\odot \epsilon
    \]
    \item decode $z$ to obtain $p_\theta(x|z)$
    \item maximize the ELBO
\end{enumerate}

\medskip

\noindent
\textbf{Generation phase.}

To sample new data:
\begin{enumerate}
    \item sample
    \[
    z \sim p(z)=\mathcal{N}(0,I)
    \]
    \item sample or decode from
    \[
    p_\theta(x|z)
    \]
\end{enumerate}

\medskip

\noindent
\textbf{Important distinction.}  
At generation time, no encoder is needed. The encoder is only used during training and inference.

\subsection{Posterior Collapse}

One of the best-known failure modes of VAEs is \textbf{posterior collapse}.

\medskip

\noindent
\textbf{Definition.}  
Posterior collapse occurs when
\[
q_\phi(z|x)\approx p(z)
\]
for many or all inputs $x$.

\medskip

\noindent
In that case, the latent variable carries little information about the input.

\medskip

\noindent
\textbf{What this means.}

\begin{itemize}
    \item the encoder effectively ignores $x$
    \item the decoder learns to reconstruct without relying on the latent code
    \item the latent representation becomes uninformative
\end{itemize}

\medskip

\noindent
\textbf{Typical causes.}

\begin{itemize}
    \item an overly expressive decoder
    \item KL regularization that is too strong too early
    \item optimization dynamics that favor ignoring $z$
\end{itemize}

\medskip

\noindent
\textbf{Intuition.}  
If the decoder is powerful enough to model the data almost on its own, then the easiest way to reduce the objective may be to keep $q_\phi(z|x)$ close to the prior and avoid using latent information.

\subsection{Practical Caveats and Common Remedies}

VAEs are elegant, but practical training requires care.

\medskip

\noindent
\textbf{Caveat 1: blurry samples.}  
Likelihood-based objectives often encourage averaging over uncertainty, which can produce samples that look smoother or blurrier than adversarial models.

\medskip

\noindent
\textbf{Caveat 2: mismatch between reconstruction and generation.}  
Good reconstructions do not automatically imply excellent unconditional samples from the prior.

\medskip

\noindent
\textbf{Caveat 3: posterior collapse.}  
As discussed above, the model may underuse the latent variables.

\medskip

\noindent
\textbf{Common remedies.}

\begin{itemize}
    \item \textbf{KL annealing:} start with a weaker KL penalty and increase it gradually
    \item \textbf{$\beta$-VAE:} replace the objective by
    \[
    \mathbb{E}_{q_\phi(z|x)}[\log p_\theta(x|z)]
    -
    \beta\,\mathrm{KL}(q_\phi(z|x)\|p(z))
    \]
    to control the tradeoff explicitly
    \item \textbf{weaker or more carefully designed decoder:} prevent the decoder from ignoring $z$
    \item \textbf{richer posterior families:} improve the approximation quality of $q_\phi(z|x)$
\end{itemize}

\medskip

\noindent
\textbf{Interpretation of $\beta$-VAE.}

\begin{itemize}
    \item $\beta > 1$ encourages stronger regularization and often more disentangled latent factors
    \item $\beta < 1$ improves reconstruction but may weaken latent structure
\end{itemize}

\subsection{Conceptual Comparison with Earlier Models}

VAEs sit at the intersection of several ideas we have seen before.

\medskip

\noindent
\textbf{From probabilistic modeling.}  
They inherit the goal of maximizing data likelihood.

\medskip

\noindent
\textbf{From latent-variable models.}  
They explain observations through hidden causes.

\medskip

\noindent
\textbf{From variational methods.}  
They replace intractable posterior inference by optimization over a tractable family.

\medskip

\noindent
\textbf{From neural networks.}  
They use flexible encoder and decoder parameterizations.

\medskip

\noindent
\textbf{Insight.}  
A VAE is not merely an autoencoder with noise; it is a deep latent-variable model trained by variational inference.

\subsection{Summary}

\textbf{Main points of this lecture.}

\begin{itemize}
    \item Generative modeling aims to learn a distribution over data rather than only a predictive rule.
    \item A latent-variable model writes
    \[
    p_\theta(x)=\int p_\theta(x|z)p(z)\,dz.
    \]
    \item Exact marginal likelihood and exact posterior inference are usually intractable.
    \item VAEs introduce an approximate posterior $q_\phi(z|x)$ and optimize the ELBO:
    \[
    \mathcal{L}(x;\theta,\phi)
    =
    \mathbb{E}_{q_\phi(z|x)}[\log p_\theta(x|z)]
    -
    \mathrm{KL}(q_\phi(z|x)\|p(z)).
    \]
    \item The encoder performs approximate inference; the decoder performs generation.
    \item The reparameterization trick makes gradient-based optimization possible.
    \item The KL term regularizes the latent space and helps align it with the prior.
    \item Posterior collapse is a major practical issue when the latent code is ignored.
\end{itemize}

\medskip

\noindent
\textbf{Looking ahead.}  
VAEs provide a principled likelihood-based route to deep generative modeling. In the next lectures, we will contrast this with adversarial training and denoising-based generation.